\documentclass[11pt,a4paper]{article}

\usepackage[utf8]{inputenc}
\usepackage[T1]{fontenc}
\usepackage{lmodern}
\usepackage[margin=0.78in]{geometry}
\usepackage[table]{xcolor}
\usepackage{array}
\usepackage{tabularx}
\usepackage{booktabs}
\usepackage{enumitem}
\usepackage{titlesec}
\usepackage{fancyhdr}
\usepackage{lastpage}
\usepackage{hyperref}

\definecolor{TextInk}{HTML}{233142}
\definecolor{HeadingBlue}{HTML}{1D3557}
\definecolor{AccentTeal}{HTML}{2A7F92}
\definecolor{AccentGold}{HTML}{C58B4E}
\definecolor{TableStripe}{HTML}{F3F8FA}
\definecolor{SoftBlue}{HTML}{EEF6F8}
\definecolor{SoftGold}{HTML}{FCF6E8}

\hypersetup{
  colorlinks=true,
  linkcolor=HeadingBlue,
  urlcolor=AccentTeal,
  pdftitle={Geography Paper 1 SL 2025 Student Work Marking Report},
  pdfauthor={IB Geography SL Preparation}
}

\color{TextInk}
\setlength{\parskip}{0.45em}
\setlength{\parindent}{0pt}
\setlength{\tabcolsep}{5pt}
\renewcommand{\arraystretch}{1.18}
\setlist[itemize]{leftmargin=1.35em,itemsep=3pt,topsep=4pt}

\newcolumntype{Y}{>{\raggedright\arraybackslash}X}
\newcolumntype{L}[1]{>{\raggedright\arraybackslash}p{#1}}
\newcommand{\term}[1]{\textcolor{HeadingBlue}{\textbf{#1}}}
\newcommand{\score}[1]{\textcolor{AccentTeal}{\textbf{#1}}}
\newcommand{\boxnote}[3]{%
\noindent\colorbox{#1}{%
  \begin{minipage}{0.965\textwidth}
  \sffamily
  \textbf{\textcolor{HeadingBlue}{#2}} #3
  \end{minipage}
}}

\titleformat{\section}
  {\normalfont\Large\sffamily\bfseries\color{HeadingBlue}}
  {}{0pt}{}
  [\vspace{0.08em}{\color{AccentGold}\titlerule[0.8pt]}]

\titleformat{\subsection}
  {\normalfont\large\sffamily\bfseries\color{AccentTeal}}
  {}{0pt}{}

\pagestyle{fancy}
\setlength{\headheight}{14pt}
\fancyhf{}
\fancyhead[L]{\sffamily\color{AccentTeal}Teacher Marking Report}
\fancyhead[R]{\sffamily\color{HeadingBlue}May 2025 Paper 1}
\fancyfoot[C]{\sffamily\color{HeadingBlue}\thepage\ of \pageref{LastPage}}
\renewcommand{\headrulewidth}{0.4pt}

\begin{document}

\begin{center}
  {\sffamily\bfseries\color{HeadingBlue}\LARGE Geography Paper 1 SL 2025 Marking Report}\\[0.25em]
  {\sffamily\color{AccentTeal}\large Options F and G Student Work}
\end{center}

\boxnote{SoftBlue}{Basis for marking.}{This report marks
\path{plans/Geography_paper_1_SL_-2025-student-work.md} against the questions
in \path{plans/Geography_paper_1_question_booklet_SL-2025.pdf}. The supplied
file \path{plans/Geography_paper_1_HLSL_markscheme-2025.pdf} appears to be a
November 2025 markscheme and its Option F/G questions do not match this question
booklet. Therefore, the short-answer marks below are awarded from the actual
question-booklet stimulus and command terms, while the 10-mark essays are
marked using the Paper 1 markbands extracted from the supplied markscheme.
Spelling and grammar are ignored unless they affect geographical meaning.}

\section{Overall Mark}

\rowcolors{2}{TableStripe}{white}
\begin{tabularx}{\textwidth}{L{0.16\textwidth}L{0.14\textwidth}Y}
\toprule
\textbf{Question} & \textbf{Mark} & \textbf{Teacher rationale} \\
\midrule
Q11 & \score{10/10} & All source-response parts are accurate. The student
correctly identifies city skyscraper farms and the rooftop greenhouse
categories, gives a valid disadvantage of vertical farming, and provides
developed explanations for a physical change in production and a human change
in consumption. \\
Q12(a) & \score{7/10} & Reaches the 7--8 band. The answer is focused on
vector-borne disease diffusion, uses malaria throughout, and makes a relevant
relative judgment that climate/geographic conditions are more important than
economic factors. It remains at the lower end because the evidence is narrow,
there is no place-specific case-study detail, and some disease terminology is
imprecise. \\
Q13 & \score{10/10} & The numerical answers are correct, and the explanations
for ecological footprint, physical influences on residential pattern, and human
influences on residential pattern all meet the marks available. \\
Q14(a) & \score{8/10} & A secure 7--8 band essay. The Shanghai answer examines
two strategies--public transport and green space--with strengths and
weaknesses for each. It does not reach 9--10 because the discussion sometimes
conflates greenhouse gases with air pollutants, and the evaluation is not
systematic enough for the top band. \\
\midrule
\textbf{Total} & \score{35/40} & Strong performance across both options. The
student is very secure on data-response and short explanation questions. The
essay marks are lost mainly for limited breadth, limited case-study precision,
and incomplete top-band evaluation. \\
\bottomrule
\end{tabularx}
\rowcolors{2}{}{}

\section{Option F: Food and Health}

\subsection{Q11: Urban, Indoor and Vertical Food Production \score{10/10}}

\rowcolors{2}{TableStripe}{white}
\begin{tabularx}{\textwidth}{L{0.16\textwidth}L{0.12\textwidth}Y}
\toprule
\textbf{Part} & \textbf{Mark} & \textbf{Rationale} \\
\midrule
11(a)(i) & \score{1/1} & \term{City skyscraper farms} is the production type
shown in all three categories: urban, indoor and vertical. \\
11(a)(ii) & \score{1/1} & \term{Urban} and \term{indoor} are the two
categories in which rooftop greenhouses are shown. \\
11(b) & \score{2/2} & The student identifies a valid disadvantage for farmers:
vertical farming limits the range of crops, especially tall-growing crops. The
answer develops this by explaining that stacked growing spaces limit crop
height. \\
11(c)(i) & \score{3/3} & The response identifies climate change/drought as a
physical influence, explains that drought risk changes over time, and develops
the production response through drought-resistant crops or irrigation. \\
11(c)(ii) & \score{3/3} & The response identifies urbanization/employment
change as a human factor, explains the shift from manual labour to office work,
and links this to demand for convenient ultra-processed foods rather than
staple-based diets. \\
\bottomrule
\end{tabularx}
\rowcolors{2}{}{}

\subsection{Q12(a): Vector-Borne Disease Diffusion \score{7/10}}

This answer fits the 7--8 markband because it addresses the whole question and
makes a clear relative judgment. It argues that geographic features and climate
are the most important controls on malaria diffusion because the mosquito
vector depends on warm conditions, standing water and suitable environments.
The answer then considers economic factors, especially health infrastructure,
vaccination/treatment availability and sanitation, as secondary influences on
the rate of diffusion.

The mark is held at \score{7/10} rather than higher because the support is
general. A top-band response would use more precise place evidence, for example
named regions where malaria transmission is limited by altitude, temperature,
drainage or public-health infrastructure. The answer also contains some
terminology weaknesses: malaria is caused by a parasite rather than a virus, and
the discussion of vaccines/treatments is relevant to control but less directly
linked to geographic diffusion than vector habitat, climate, water and
mobility.

To move this to \score{9--10}, the student should compare several geographic
factors more systematically: climate, altitude, water availability, land use,
settlement density, mobility, poverty and public-health infrastructure. The
final judgment should explain where and why one set of factors becomes more
important than another.

\section{Option G: Urban Environments}

\subsection{Q13: The 15-Minute City \score{10/10}}

\rowcolors{2}{TableStripe}{white}
\begin{tabularx}{\textwidth}{L{0.16\textwidth}L{0.12\textwidth}Y}
\toprule
\textbf{Part} & \textbf{Mark} & \textbf{Rationale} \\
\midrule
13(a)(i) & \score{1/1} & \term{13.5\%} is an acceptable estimate for citizens
more than 15 minutes from essential destinations. \\
13(a)(ii) & \score{1/1} & \term{6 minutes} is correct: the six times
4, 5, 4, 6.5, 8.5 and 8 minutes average to 6 minutes. \\
13(b) & \score{2/2} & The response states that a 15-minute city reduces
dependence on cars and develops this with lower greenhouse-gas emissions. This
is a valid way to reduce ecological footprint. \\
13(c)(i) & \score{3/3} & Mountains/steep slopes are a valid physical factor.
The answer explains that building on steep slopes is harder, so residential
areas are more likely to be built around them rather than on them. \\
13(c)(ii) & \score{3/3} & Access to essential goods and services is a valid
human factor. The answer explains that residences are more likely to cluster
near services, especially where congestion makes car travel less attractive and
walkability more important. \\
\bottomrule
\end{tabularx}
\rowcolors{2}{}{}

\subsection{Q14(a): Reducing Air Pollution in Shanghai \score{8/10}}

This essay securely reaches the 7--8 band. It is focused on the command term
\term{examine} and on the required issue of strategies to reduce air pollution
in an urban area. Shanghai is used consistently as the urban example. The
student examines two strategies: improving public transportation and increasing
green space.

The public transport paragraph has clear strengths: high-capacity transit can
replace many individual car journeys, reduce traffic congestion and reduce
idling emissions. It also gives useful case-specific support, including the
long metro network and more than 10 million daily riders. The weakness is also
valid: public transport has high construction costs and may still rely on
electricity generated from fossil fuels.

The green-space paragraph is also balanced. The student explains that trees can
trap particulate matter and absorb pollutants such as $CO_2$, $NO_2$ and
$SO_2$, and can cool urban areas through shade and evapotranspiration. The
weakness is relevant: green space has a weaker direct impact than reducing
anthropogenic emissions and is seasonally dependent where trees lose leaves.

The essay does not reach \score{9--10} because the evaluation is not fully
systematic. Some discussion is stronger for greenhouse-gas reduction than for
local air pollution, and the answer would be sharper if it separated
particulate matter, nitrogen oxides, sulphur dioxide and carbon dioxide. A
top-band answer would compare the two strategies directly, explain which is
more effective for different pollutants and time scales, and support the
judgment with more precise Shanghai evidence.

\section{Main Targets}

\begin{itemize}
  \item In 10-mark essays, add named places and statistics that directly
  support the argument.
  \item Separate air pollutants from greenhouse gases when the question is
  specifically about air pollution.
  \item Make the final judgment more comparative: state which factor or
  strategy is more important, for which places, pollutants or time scales, and
  why.
  \item Keep the clear short-answer structure; the data-response and 3-mark
  explanations are already very strong.
\end{itemize}

\end{document}
